\section{Lyapunov Based Control Design}
\subsection{Control Lyapunov Functions}\label{CLF}
Consider the stabilization problem for the following system
\begin{equation*}
    \dot{x} = f(x) + \sum_{i=1}^m g_i(x) u_i = f(x) + G(x)u,
\end{equation*}
$f:\mathbb{R}^n\to\mathbb{R}^n$, $g_i : \mathbb{R}^n \to \mathbb{R}^n$, $G : \mathbb{R}^n \to \mathbb{R}^{n \times m}$, $u : \mathbb{R} \to \mathbb{R}^m$.
\newpar{}
A function $V : \mathbb{R}^n \to \mathbb{R}$ is a control Lyapunov function (CLF) for the system above if it is differentiable, positive definite, radially unbounded, and
\begin{equation*}
    \forall x \neq 0, \quad \inf_u \frac{\partial V}{\partial x}(f(x) + G(x)u) < 0
\end{equation*}
For physical systems: possibility to choose $u$ such that the energy decreases.

\subsubsection{Conditions to Find Control Lyapunov Functions}

The condition~\ref{CLF} shown above can be hard (or impossible) to check since it involves solving an optimization problem for each $x \neq 0$. A more practical way is the following:
\newpar{}
Condition above is equivalent to
\begin{equation*}
    \forall x \neq 0: \quad \frac{\partial V}{\partial x} G(x) = 0 \Rightarrow \frac{\partial V}{\partial x} f(x) < 0
\end{equation*}
i.e.\ if we lose control at some point $x$, then the system must autonomously converge, i.e.\ $\frac{\partial V}{\partial x} f(x)<0$.
\newpar{}
The condition can also be reformulated as:
\begin{equation*}
    \forall x \neq 0: \quad \frac{\partial V}{\partial x} f(x) \geq 0 \Rightarrow \frac{\partial V}{\partial x} G(x) \neq 0,
\end{equation*}
i.e.\ for points where the system tends to diverge, we must have control in the direction of $\dot{V}$.

\subsection{Backstepping}
In this subsection, we show how a CLF can be constructed for the special case of systems in strict-feedback form, by using backstepping. It is a recursive approach that constructs a sequence of (simpler) virtual subsystems, each of which is stabilized individually.
\subsubsection{Integrator backstepping}\label{int_backstepping}
Consider the following system:
\begin{align*}
    \dot{x}_1 & = f_1(x_1) + g_1(x_1)x_2 \\
    \dot{x}_2 & = u
\end{align*}
where $x_1 \in \mathbb{R}^{n_1}$, $x_2, u \in \mathbb{R}$ (single input).
\newpar{}
We want to stabilize the origin of $\dot{x_1}$. The system can be viewed as a cascade of two subsystems: one with $x_2$ as input, and the second with $u$ as input.
\newpar{}
\ptitle{Assumption}

We know a (smooth) state-FB control $\alpha_1 : \mathbb{R}^{n_1} \to \mathbb{R}$ and a smooth, positive definite CLF $V_1 : \mathbb{R}^{n_1} \to \mathbb{R}$ such that:
\begin{equation*}
    \dot{V}_1 = \frac{\partial V_1}{\partial x_1} \left[f_1(x_1) + g_1(x_1)\alpha_1(x_1)\right] \leq -W_1(x_1),
\end{equation*}
where $W_1(x)$ is a positive definite function, i.e.\ the origin of $x_1$ is asymptotically stable.

\paragraph{Control Approach}
\ptitle{Error Coordinates}

Using the reformulation
\begin{align*}
    \dot{x}_1 & = f_1(x_1) + g_1(x_1)\alpha_1(x_1) + g_1(x_1)[x_2 - \alpha_1(x_1)] \\
    \dot{x}_2 & = u
\end{align*}
we can introduce the \textbf{error coordinates}
\begin{equation*}
    (e_1, e_2) := (x_1 - 0, x_2 - \alpha_1(x_1))
\end{equation*}
\begin{align*}
    \dot{e}_1 & = f_1(e_1) + g_1(e_1)\alpha_1(e_1) + g_1(e_1)e_2                                 \\
    \dot{e}_2 & = u - \dot{\alpha}_1 = u - \frac{\partial \alpha_1}{\partial e_1} \dot{e}_1 =: v
\end{align*}
Now, find $u = \alpha_2(x_1, x_2)$ such that $(e_1, e_2) \to 0$ as $t \to \infty$. In that case, we have for $\dot{x}_1$ the dynamics according to above reformulation which is asymptotically stable by assumption.

\newpar{}
\ptitle{Lyapunov Function}

Define $V_2(e_2) := \frac{1}{2}e_2^2$ and $V(e_1, e_2) := V_1(e_1) + V_2(e_2)$. Then:
{\small
\begin{gather*}
    \dot{V}(e_1, e_2) \\
    = \frac{\partial V_1}{\partial e_1}(f_1(e_1) + g_1(e_1)\alpha_1(e_1)) + \frac{\partial V_1}{\partial e_1}g_1(e_1)e_2 + e_2u - e_2\dot{\alpha}_1\\
    \leq -W_1(e_1) + \frac{\partial V_1}{\partial e_1}g_1(e_1)e_2 + e_2u - e_2\dot{\alpha}_1
\end{gather*}
}
\newpar{}
\ptitle{Control Input}\label{backstepping_input}

To make $\dot{V}$ negative definite, choose
\begin{equation*}
    u := \underbrace{-\frac{\partial V_1}{\partial e_1}g_1(e_1) + \dot{\alpha}_1}_{\text{cancelling}} \underbrace{- k_2e_2}_{\text{stabilization}} =: \alpha_2(e_1, e_2), \quad k_2 > 0
\end{equation*}
which yields
\begin{equation*}
    \dot{V}(e_1, e_2) = -W_1(e_1) - k_2 e_2^2 < 0, \quad \forall (e_1, e_2) \neq (0, 0)
\end{equation*}
i.e.\ $(e_1, e_2)$ is asymptotically stable.
\newpar{}
\ptitle{Convergence}

This result does not necessarily imply that $(x_1, x_2) = (0, 0)$ is an asymptotically stable equilibrium for the original system. In particular, in original coordinates we have
\begin{equation*}
    u = -\frac{\partial V_1}{\partial x_1}g_1(x_1) + \dot{\alpha}_1 - k_2(x_2 - \alpha_1(x_1)) = \alpha_2(x_1, x_2),
\end{equation*}
for which we see that to ensure $(x_1, x_2) = (0, 0)$ is asymptotically stable, we need
\begin{equation*}
    \alpha_1(0) = 0
\end{equation*}

\subsubsection{First Generalization}
Let us now generalize by including functions $f_2 : \mathbb{R} \times \mathbb{R} \to \mathbb{R}$ and $g_2 : \mathbb{R} \times \mathbb{R} \to \mathbb{R}$:
\begin{align*}
    \dot{x}_1 & = f_1(x_1) + g_1(x_1)x_2         \\
    \dot{x}_2 & = f_2(x_1, x_2) + g_2(x_1, x_2)u
\end{align*}
with $x_2, u_2 \in \mathbb{R}, g_2(x_1, x_2) \neq 0, \forall x_1, x_2$ and $g_2$ is known in sign and intensity.
\paragraph{Control Approach}
\ptitle{Virtual Control Input}

Using the transformation
\begin{equation*}
    u = \frac{1}{g_2(x_1, x_2)} (v - f_2(x_1, x_2)),
\end{equation*}
we get:
\begin{align*}
    \dot{x}_1 & = f_1(x_1) + g_1(x_1)x_2 \\
    \dot{x}_2 & = v
\end{align*}
which matches the structure from~\ref{int_backstepping}.

Therefore, by following the same steps, we define $v$ to render $(x_1, x_2) = (0, 0)$ asymptotically stable:
{\footnotesize
\begin{equation*}
    u = \frac{1}{g_2(x_1, x_2)}\left(\underbrace{-\frac{\partial V_1}{\partial x_1}g_1(x_1) + \dot{\alpha}_1 - k_2(x_2 - \alpha_1(x_1))}_{=:v} - f_2(x_1, x_2)\right)
\end{equation*}
}
Note that $v$ is computed as as shown in~\ref{backstepping_input}.

\subsubsection{Generalization to Three States}

In the following we see backstepping through two integrators:   % TODO: Integrators correct?
\begin{align*}
    \dot{x}_1 & = f_1(x_1) + g_1(x_1)x_2                   \\
    \dot{x}_2 & = f_2(x_1, x_2) + g_2(x_1, x_2)x_3         \\
    \dot{x}_3 & = f_3(x_1, x_2, x_3) + g_3(x_1, x_2, x_3)u
\end{align*}
Assume: $g_2, g_3 \ne 0, \, \forall x_1, x_2, x_3$, and both their value and their sign are known.
\paragraph{Control Approach}
\ptitle{Error Coordinates}

In the part above, we showed that there exists $\alpha_2$ such that for $x_3 = \alpha_2(x_1, x_2)$, $(e_1, e_2) \to 0$. Now, we want to define the input $u$ so that $x_3$ approaches $\alpha_2$, and simultaneously $x_3 \to 0$. Thus, we introduce the error coordinate $e_3 := x_3 - \alpha_2(x_1, x_2)$.

Then, we can use the following input transformation:
\begin{equation*}
    u = \frac{1}{g_3(x_1, x_2, x_3)}\left(v - f_3(x_1, x_2, x_3)\right),
\end{equation*}
which simplifies the system as follows:
\begin{align*}
    \dot{x}_1 & = f_1(x_1) + g_1(x_1)x_2           \\
    \dot{x}_2 & = f_2(x_1, x_2) + g_2(x_1, x_2)x_3 \\
    \dot{x}_3 & = v,
\end{align*}
and therefore we can apply integrator backstepping once more.

In "error coordinates", the system reads:
\begin{align*}
    \dot{e}_1 & = f_1(e_1) + g_1(e_1)(e_2 + \alpha_1(e_1))                                                 \\
    \dot{e}_2 & = f_2(e_1, e_2 + \alpha_1(e_1))                                                            \\
              & + g_2(e_1, e_2 + \alpha_1(e_1))(e_3 + \alpha_2(e_1, e_2 + \alpha_1(e_1))) - \dot{\alpha}_1 \\
    \dot{e}_3 & = v - \dot{\alpha}_2
\end{align*}
\newpar{}
\ptitle{Lyapunov Function}

Now, define $V_3(e_3) := \frac{1}{2}e_3^2$, and $V(e_1, e_2, e_3) := V_1(e_1) + V_2(e_2) + V_3(e_3)$. One can show that:
\begin{equation*}
    \dot{V} = L_{f_1 + g_1\alpha_1}V_1(e_1) - k_2e_2^2 + e_2g_2(e_1, e_2 + \alpha_1(e_1))e_3 + e_3(v - \dot{\alpha}_2)
\end{equation*}

Since we want to ensure that $\dot{V}$ is negative definite, then we can choose the variable $v$ based on the same reasoning discussed above, i.e., with some "cancelling terms", and some "stabilizing terms":
\begin{equation*}
    v = -\underbrace{e_2g_2(e_1, e_2 + \alpha_1(e_1)) + \dot{\alpha}_2}_{\text{``cancelling terms''}} - \underbrace{k_3e_3}_{\text{``stabilizing term''}}.
\end{equation*}
\newpar{}
\ptitle{Control Input}

Then, in original coordinates:
{\tiny
\begin{equation*}
    u  = \frac{- (x_2 - \alpha_1(x_1))g_2(x_1, x_2) + \dot{\alpha}_2 - k_3(x_3 - \alpha_2(x_1, x_2)) - f_3(x_1, x_2, x_3)}{g_3(x_1, x_2, x_3)}
\end{equation*}
}
\begin{equation*}
    =: \alpha_3(x_1, x_2, x_3)
\end{equation*}

\subsubsection{General backstepping recursion}

Based on the two generalizations shown above, it is now possible to employ the backstepping technique to a general system in \textbf{strict-feedback form}:
\begin{align*}
    \dot{x}_1 & = f_1(x_1) + g_1(x_1)x_2                         \\
    \dot{x}_2 & = f_2(x_1, x_2) + g_2(x_1, x_2)x_3               \\
              & \vdots                                           \\
    \dot{x}_k & = f_k(x_1, \ldots, x_k) + g_k(x_1, \ldots, x_k)u
\end{align*}

$x_1 \in \mathbb{R}^{n_1}, \, x_2, \ldots, x_k, u \in \mathbb{R}, \, g_2, \ldots, g_k \text{ nowhere zero}.$

The functions $f_i, g_i$ are (sufficiently) smooth, because their derivatives are needed in the definitions of $\alpha_i$.

\newpar{}
\ptitle{Backstepping Recursion}

\begin{enumerate}
    \item Input data: A CLF $V_1$ for the first system $\dot{x}_1 = f_1(x_1) + g_1(x_1)u_1$, with a smooth feedback $\alpha_1$ that asymptotically stabilizes the origin for $u_1 = \alpha_1(x_1)$.
    \item For $i = 2, \ldots, k$:
          Construct a CLF $V_i(e_i) = \frac{1}{2}e_i^2$, define $V = \sum_{j=1}^i V_j$, and a feedback $\alpha_i$ that asymptotically stabilizes
          $(e_1, \ldots, e_i) = (x_1 - 0, x_2 - \alpha_1(x_1), \ldots, x_i - \alpha_{i-1}(x_1, \ldots, x_{i-1}))$
          with
              {\footnotesize
                  \begin{gather*}
                      \alpha_i(x_1, \ldots, x_i) \\
                      = \frac{1}{g_i}\left( -\frac{\partial V_{i-1}}{\partial e_{i-1}}g_{i-1} + \dot{\alpha}_{i-1} - k_i(x_i - \alpha_{i-1}(x_1, \ldots, x_{i-1})) - f_i \right)
                  \end{gather*}
              }
    \item Finally, the input is chosen as:
          \begin{equation*}
              u := \alpha_k(x_1, \ldots, x_k)
          \end{equation*}
\end{enumerate}
\textbf{Note} that the function $V = \sum_{i=1}^k V_i$ with $V_i(e_i) = \frac{1}{2}e_i^2$ for $i = 2, \ldots, k$, is a Control Lyapunov Function.

\paragraph{Disturbance Rejection}
In the part above, we only focused on the case where the system is not affected by disturbances. However, it is important to remark that backstepping can also be used with systems affected by \textbf{bounded disturbances}.
